\HuazhongGuide{本表参照以下提纲撰写，要求逻辑清晰，层次分明，内容翔实，排版规范。（不超过7000 字）}

\section{\textbf{项目的立项依据。}研究目的、意义和国内外概况（附主要参考文献）}

% 迁移来源：extraTex/1.1.立项依据.tex；全文（含原始插图与交叉引用）
\subsection{研究背景与意义}

本项目拟研究高等教育场景下教育智能体的可信评测与伦理风险监测机制。大语言模型与检索、工具调用和记忆模块结合后，能够连续完成课程资料查找、学习提示生成及文献整理等任务。ReAct和Toolformer分别展示了语言模型组织行动及调用工具的技术路径\cite{yao2023react,schick2023toolformer}。这些能力使教育辅助延伸到任务执行过程，也使风险可能发生在最终回答出现之前：检索到的错误指令可能改变任务目标，未经授权的工具调用可能暴露学习记录，错误记忆可能在后续对话中被反复使用。仅审查最终答案，无法充分描述这些已经发生的行为。

教育情境还对“正确”的含义提出了额外要求。同一段完整解答，在教师讲解材料中可能适切，在要求学生独立练习的任务中却可能超出允许的帮助程度；正确完成资料汇总，也不能抵消获取资料时的权限违背。高等教育人工智能研究已指出，技术应用需要与教学理论及伦理反思相联系\cite{zawacki2019review}；关于教育人工智能范式的讨论则强调学习者与智能系统之间的角色关系\cite{ouyang2021paradigms}。因此，教育智能体的可信性需要在任务目的、支持边界与行动授权共同构成的情境中判断。

已有教育试验提示，辅助过程的设计可能改变结果的方向。Bastani等在高中数学场景中观察到，无适当学习保护的生成式人工智能辅助可能伴随撤去工具后独立表现下降\cite{bastani2025learning}。Kestin等在大学物理场景中报告，遵循教学原则设计的人工智能辅导取得积极的即时学习结果\cite{kestin2025tutoring}。两项研究的对象和设计不同，不能据此直接判断某一高校智能体的效果，但共同说明使用过程值得单独考察。面向教学与人机交互的综合讨论也将模型局限、学习者能动性和教育应用设计置于同一视野\cite{yan2024learning}。

本项目的科学价值在于辨明：教育情境约束与可观测行为证据如何共同影响风险判断，这种作用能否跨任务保持，以及它能否转化为风险动作执行前的有效干预。拟形成的评测框架可为教育工具选型、案例复核和运行监测提供可追溯依据。研究按一年期开放课题组织，先在公开资料与可控沙箱上建立证据，再在条件落实后开展小规模高校案例验证。
见图~\ref{fig:edu-f01}。
\HuazhongFigure{figures/F01-background-v2.pdf}{教育情境与过程证据共同决定风险含义}{fig:edu-f01}



\subsection{国内外研究现状及发展动态分析}

\textbf{教育评价已从回答质量延伸到教学行为，但不同终点仍需区分。}Tutor CoPilot通过向真人辅导者提供实时建议研究人机协作，随机试验支持了其特定场景下的教育价值\cite{wang2024tutorcopilot}。此类教学支持研究仍需区分不同评价终点：专家偏好、短时任务完成与长期学习保持不能相互替代。本项目据此将教学支持是否符合任务要求作为过程评价对象，避免以模型评分替代学习成效。

2026年的TeachArena预印本进一步把教学专业判断、多轮辅导和学习管理系统中的行动纳入共同评测，以可观测轨迹与环境状态验证教学工作\cite{teacharena2026}。这一进展说明教育智能体过程评价已有直接研究基础。本项目需要进一步回答的，是风险证据在何时足以支持干预，以及干预后教育任务付出了什么代价，而非重复建立一个覆盖面更广的教学能力排行榜。

\textbf{通用智能体基准已覆盖行为与环境后果，教育约束仍需专门操作化。}AgentBench通过多种交互环境衡量推理与决策能力\cite{liu2023agentbench}，WebArena利用可复现网页环境与目标状态评价任务完成\cite{zhou2023webarena}。ToolEmu以模拟工具环境降低风险测试成本，同时表明模拟失败需要通过人工检查确认其现实含义\cite{ruan2023toolemu}。R-Judge直接围绕多轮交互记录进行风险识别\cite{yuan2024rjudge}。这些工作已将评价对象从最终答案拓展至行为轨迹及环境后果。

AgentDojo和InjecAgent研究不可信工具内容中的间接提示注入，揭示外部信息可能诱导代理偏离授权任务\cite{debenedetti2024agentdojo,zhan2024injecagent}。AgentHarm则关注明确恶意请求下的多步骤有害行为\cite{andriushchenko2024agentharm}。不同基准的威胁模型各有边界；对教育任务而言，正常学习请求中的证据误用、帮助程度失配与权限越界，需要结合具体教学目的辨认。公开恶意任务上的拒绝能力不能直接代表正常学习过程中发现隐性风险的能力。

\textbf{过程评判提供了可借鉴的方法，评价器本身也需要验证。}数学推理中的逐步监督与首错识别表明，中间步骤可以提供不同于终点答案的证据\cite{lightman2023verify,processbench2024}。但工具执行记录中的风险未必对应一个逻辑错误：某项操作可能在技术上完全成功，却违反本次任务的授权范围。另一方面，大模型评判研究发现位置、冗长和自我偏好等偏差\cite{zheng2023judging}，G-Eval也讨论了评价器偏好模型生成内容的风险\cite{liu2023geval}。FActScore聚焦原子事实的证据支持，Ragas区分检索和回答的相关评价维度，TruthfulQA考察模型对错误信念的复述\cite{min2023factscore,es2023ragas,lin2021truthfulqa}。这些指标可以帮助设计局部判据，但其相加不能自动形成具有测量效度的教育可信总分。

\textbf{运行保护已形成多种技术路径，安全与任务效用的比较仍是关键。}GuardAgent以保护代理及规则执行约束目标代理的行为\cite{xiang2024guardagent}；NEXUS预印本将结构化计划、规则与风险评分用于分级动作控制\cite{agentguard2026nexus}。因此，在线监测与阻断本身已有研究基础，进一步的价值应体现在对特定风险的有效性及其代价上。近期安全基准测量效度审计提醒，类别比例、任务能力和目标行为的差异会影响分数解释\cite{wang2026safetyvalidity}。该工作仍属预印本，本项目将其作为方法学警示，采用配对任务、独立校准和安全与效用分开报告的设计。

\subsection{现有研究的不足}

首先，教育情境约束与轨迹证据之间的关系尚需在可控条件下检验。教学目标、允许提供的帮助和数据访问范围往往同时变化，评价器可能只学到显式风险词或任务模板，而没有利用情境信息。若仅比较“有轨迹”与“无轨迹”，提升还可能来自信息量或文本长度差异，不能据此证明教育约束产生了独立作用。需要设置同等可观测信息的通用轨迹评价器，并通过情境约束与证据消融识别增益来源。

其次，完整轨迹的事后识别与运行中的提前告警之间存在时间差。事后评价能利用最终答案与执行后果，运行监测只能访问当前已经产生的信息。如果用完整轨迹标签或未来动作辅助当前评分，实验会高估提前发现能力。教育场景还需要区分首次可观察风险、首次告警和风险动作提交三个时点，使“及时”具有可复核含义。

最后，风险减少与教学支持受损需要在同一组初始任务上共同观察。频繁拒绝或移交人工可以减少风险动作，但也可能使本应完成的学习任务中断。仅报告拦截率无法判断干预是否值得采用。真实高校案例又受到课程、授权和样本规模限制，沙箱有效性不能直接外推为普遍教学效果。由此需要把受控机制检验与小规模案例验证分开，并公开其解释边界。

\subsection{关键科学问题与研究切入点}

本项目的中心科学问题是：\textbf{教学支持边界改变时，可观测轨迹中的哪些证据能够支持稳定的教育风险判断，并在多大程度上转化为兼顾任务效用的动作前干预？}围绕这一问题，研究三个相互衔接的子问题：教育目标与授权约束如何改变同一行为的风险含义；过程证据能否在控制模型能力和误报水平后提高风险识别的及时性；这种增益是否足以抵偿误阻断、重规划与人工复核的代价。

据此提出科学假说：在任务家族独立、模型与调用预算一致的条件下，正确教育约束能够增强轨迹证据对支持边界风险的识别贡献，该增益大于约束打乱时的相应增益。进一步假设，当有效风险证据早于动作提交出现时，基于该证据的分级干预能够在正常任务自主完成率损失受限的条件下降低风险执行比例。

研究通过教育约束与轨迹证据交叉消融、统一的名义误报校准、家族留出及干预对照进行检验，区分上述作用与风险词、模板和额外拒绝的影响。效应区间下界超过最小实际意义幅度时支持预定增益，上界低于该幅度时不支持达到预定增益，跨界时保留不确定判断；干预收益需同时获得风险执行下降与任务效用损失满足约束的区间支持。家族差异用于界定适用范围，实际误报差异纳入结果解释。

研究将以“风险表征—轨迹评测—监测干预”三条路线展开。第一条路线确定风险事件及其可观察证据，第二条路线检验不同证据的识别作用与时间边界，第三条路线在相同任务中检验干预后果。课程答疑和文献学习辅助作为共用验证场景，双智能体传播作为受控压力测试。项目预期贡献是给出有适用条件的测量方法、风险案例和监测原型，而不是把未经验证的单一得分作为教育智能体可信性的充分证明。

% 迁移来源：references/myexample.bib；由原始题录生成的全部被引文献；无模板样式设置
% 由 myexample.bib 和仓库共享 gbt7714-nsfc.bst 生成；28条被引文献。
% 原第9条 learnlm2024 已按用户要求删除；其余条目保持原相对顺序。
% 重生成：python3 tests/ccnu-eduagenttrust-template/20260919-references/render_references.py
\begin{thebibliography}{10}
\providecommand{\bibauthor}[1]{#1}
\providecommand{\bibeditor}[1]{#1}
\providecommand{\bibtranslator}[1]{#1}
\providecommand{\bibtitle}[1]{#1}
\providecommand{\bibbooktitle}[1]{#1}
\providecommand{\bibjournal}[1]{#1}
\providecommand{\bibmark}[1]{#1}
\providecommand{\bibcountry}[1]{#1}
\providecommand{\bibpatentid}[1]{#1}
\providecommand{\bibedition}[1]{#1}
\providecommand{\biborganization}[1]{#1}
\providecommand{\bibaddress}[1]{#1}
\providecommand{\bibpublisher}[1]{#1}
\providecommand{\bibinstitution}[1]{#1}
\providecommand{\bibschool}[1]{#1}
\providecommand{\bibvolume}[1]{#1}
\providecommand{\bibnumber}[1]{#1}
\providecommand{\bibversion}[1]{#1}
\providecommand{\bibpages}[1]{#1}
\providecommand{\bibmodifydate}[1]{#1}
\providecommand{\bibcitedate}[1]{#1}
\providecommand{\bibyear}[1]{#1}
\providecommand{\bibdate}[1]{#1}
\providecommand{\biburl}[1]{\newline\url{#1}}

\bibitem{yao2023react}
\bibauthor{Yao S\@, Zhao J\@, Yu D\@, Du N\@, Shafran I\@, Narasimhan K\@, Cao
  Y}\@. \bibtitle{{ReAct: Synergizing Reasoning and Acting in Language
  Models}}\bibmark{[J/OL]}\@. \bibjournal{arXiv preprint arXiv:2210.03629}\@,
  \bibyear{2022}\bibcitedate{[2026-09-06]}\@.
  \url{https://arxiv.org/abs/2210.03629}\@.

\bibitem{schick2023toolformer}
\bibauthor{Schick T\@, Dwivedi-Yu J\@, Dessì R\@, Raileanu R\@, Lomeli M\@,
  Zettlemoyer L\@, Cancedda N\@, Scialom T}\@. \bibtitle{{Toolformer: Language
  Models Can Teach Themselves to Use Tools}}\bibmark{[J/OL]}\@.
  \bibjournal{arXiv preprint arXiv:2302.04761}\@,
  \bibyear{2023}\bibcitedate{[2026-09-06]}\@.
  \url{https://arxiv.org/abs/2302.04761}\@.

\bibitem{zawacki2019review}
\bibauthor{Zawacki-Richter O\@, Marín V~I\@, Bond M\@, Gouverneur F}\@.
  \bibtitle{{Systematic review of research on artificial intelligence
  applications in higher education -- where are the
  educators?}}\bibmark{[J/OL]}\@. \bibjournal{International Journal of
  Educational Technology in Higher Education}\@,
  \bibvolume{16}\bibnumber{(1)}\thinspace{}\textnormal{: }\bibpages{39}\@,
  \bibyear{2019}\bibcitedate{[2026-09-06]}\@.
  \url{https://doi.org/10.1186/s41239-019-0171-0}\@.

\bibitem{ouyang2021paradigms}
\bibauthor{Ouyang F\@, Jiao P}\@. \bibtitle{{Artificial intelligence in
  education: The three paradigms}}\bibmark{[J/OL]}\@. \bibjournal{Computers and
  Education: Artificial Intelligence}\@, \bibvolume{2}\thinspace{}\textnormal{:
  }\bibpages{100020}\@, \bibyear{2021}\bibcitedate{[2026-09-06]}\@.
  \url{https://doi.org/10.1016/j.caeai.2021.100020}\@.

\bibitem{bastani2025learning}
\bibauthor{Bastani H\@, Bastani O\@, Sungu A\@, Ge H\@, Özge Kabakcı\@,
  Mariman R}\@. \bibtitle{{Generative AI without guardrails can harm learning:
  Evidence from high school mathematics}}\bibmark{[J/OL]}\@.
  \bibjournal{Proceedings of the National Academy of Sciences}\@,
  \bibvolume{122}\bibnumber{(26)}\thinspace{}\textnormal{:
  }\bibpages{e2422633122}\@, \bibyear{2025}\bibcitedate{[2026-09-06]}\@.
  \url{https://doi.org/10.1073/pnas.2422633122}\@.

\bibitem{kestin2025tutoring}
\bibauthor{Kestin G\@, Miller K\@, Klales A\@, Milbourne T\@, Ponti G}\@.
  \bibtitle{{AI tutoring outperforms in-class active learning: an RCT
  introducing a novel research-based design in an authentic educational
  setting}}\bibmark{[J/OL]}\@. \bibjournal{Scientific Reports}\@,
  \bibvolume{15}\bibnumber{(1)}\thinspace{}\textnormal{: }\bibpages{17458}\@,
  \bibyear{2025}\bibcitedate{[2026-09-06]}\@.
  \url{https://doi.org/10.1038/s41598-025-97652-6}\@.

\bibitem{yan2024learning}
\bibauthor{Yan L\@, Greiff S\@, Teuber Z\@, Gašević D}\@. \bibtitle{{Promises
  and challenges of generative artificial intelligence for human
  learning}}\bibmark{[J/OL]}\@. \bibjournal{Nature Human Behaviour}\@,
  \bibvolume{8}\bibnumber{(10)}\thinspace{}\textnormal{:
  }\bibpages{1839\thinspace{}\textnormal{--}\thinspace{}1850}\@,
  \bibyear{2024}\bibcitedate{[2026-09-06]}\@.
  \url{https://doi.org/10.1038/s41562-024-02004-5}\@.

\bibitem{wang2024tutorcopilot}
\bibauthor{Wang R~E\@, Ribeiro A~T\@, Robinson C~D\@, Loeb S\@, Demszky D}\@.
  \bibtitle{{Tutor CoPilot: A Human-AI Approach for Scaling Real-Time
  Expertise}}\bibmark{[J/OL]}\@. \bibjournal{arXiv preprint
  arXiv:2410.03017}\@, \bibyear{2024}\bibcitedate{[2026-09-06]}\@.
  \url{https://arxiv.org/abs/2410.03017}\@.

\bibitem{teacharena2026}
\bibauthor{Chen Z\@, Liu P\@, Sheng R\@, Li H\@, Tu J\@, Deng X\@, Shum K\@,
  Liu D\@, Qu H}\@. \bibtitle{{TeachArena: Are Language Agents Ready for
  Realistic Teaching Work?}}\bibmark{[J/OL]}\@. \bibjournal{arXiv preprint
  arXiv:2605.14322}\@, \bibyear{2026}\bibcitedate{[2026-09-06]}\@.
  \url{https://arxiv.org/abs/2605.14322}\@.

\bibitem{liu2023agentbench}
\bibauthor{Liu X\@, Yu H\@, Zhang H\@, Xu Y\@, Lei X\@, Lai H\@, Gu Y\@, Ding
  H\@, Men K\@, Yang K\@, Zhang S\@, Deng X\@, Zeng A\@, Du Z\@, Zhang C\@,
  Shen S\@, Zhang T\@, Su Y\@, Sun H\@, Huang M\@, Dong Y\@, Tang J}\@.
  \bibtitle{{AgentBench: Evaluating LLMs as Agents}}\bibmark{[J/OL]}\@.
  \bibjournal{arXiv preprint arXiv:2308.03688}\@,
  \bibyear{2023}\bibcitedate{[2026-09-06]}\@.
  \url{https://arxiv.org/abs/2308.03688}\@.

\bibitem{zhou2023webarena}
\bibauthor{Zhou S\@, Xu F~F\@, Zhu H\@, Zhou X\@, Lo R\@, Sridhar A\@, Cheng
  X\@, Ou T\@, Bisk Y\@, Fried D\@, Alon U\@, Neubig G}\@. \bibtitle{{WebArena:
  A Realistic Web Environment for Building Autonomous
  Agents}}\bibmark{[J/OL]}\@. \bibjournal{arXiv preprint arXiv:2307.13854}\@,
  \bibyear{2023}\bibcitedate{[2026-09-06]}\@.
  \url{https://arxiv.org/abs/2307.13854}\@.

\bibitem{ruan2023toolemu}
\bibauthor{Ruan Y\@, Dong H\@, Wang A\@, Pitis S\@, Zhou Y\@, Ba J\@, Dubois
  Y\@, Maddison C~J\@, Hashimoto T}\@. \bibtitle{{Identifying the Risks of LM
  Agents with an LM-Emulated Sandbox}}\bibmark{[J/OL]}\@. \bibjournal{arXiv
  preprint arXiv:2309.15817}\@, \bibyear{2023}\bibcitedate{[2026-09-06]}\@.
  \url{https://arxiv.org/abs/2309.15817}\@.

\bibitem{yuan2024rjudge}
\bibauthor{Yuan T\@, He Z\@, Dong L\@, Wang Y\@, Zhao R\@, Xia T\@, Xu L\@,
  Zhou B\@, Li F\@, Zhang Z\@, Wang R\@, Liu G}\@. \bibtitle{{R-Judge:
  Benchmarking Safety Risk Awareness for LLM Agents}}\bibmark{[J/OL]}\@.
  \bibjournal{arXiv preprint arXiv:2401.10019}\@,
  \bibyear{2024}\bibcitedate{[2026-09-06]}\@.
  \url{https://arxiv.org/abs/2401.10019}\@.

\bibitem{debenedetti2024agentdojo}
\bibauthor{Debenedetti E\@, Zhang J\@, Balunović M\@, Beurer-Kellner L\@,
  Fischer M\@, Tramèr F}\@. \bibtitle{{AgentDojo: A Dynamic Environment to
  Evaluate Prompt Injection Attacks and Defenses for LLM
  Agents}}\bibmark{[J/OL]}\@. \bibjournal{arXiv preprint arXiv:2406.13352}\@,
  \bibyear{2024}\bibcitedate{[2026-09-06]}\@.
  \url{https://arxiv.org/abs/2406.13352}\@.

\bibitem{zhan2024injecagent}
\bibauthor{Zhan Q\@, Liang Z\@, Ying Z\@, Kang D}\@. \bibtitle{{I}njec{A}gent:
  Benchmarking Indirect Prompt Injections in Tool-Integrated Large Language
  Model Agents}\bibmark{[C/OL]} //\thinspace{}\bibeditor{Ku L-W\@, Martins A\@,
  Srikumar V}\@. \bibbooktitle{Findings of the Association for Computational
  Linguistics: ACL 2024}\@. \bibaddress{Bangkok,
  Thailand}\thinspace{}\textnormal{: }\bibpublisher{Association for
  Computational Linguistics}\@, \bibyear{2024}\thinspace{}\textnormal{:
  }\bibpages{10471\thinspace{}\textnormal{--}\thinspace{}10506}\bibcitedate{[2026-09-06]}\@.
  \url{https://aclanthology.org/2024.findings-acl.624/}\@.

\bibitem{andriushchenko2024agentharm}
\bibauthor{Andriushchenko M\@, Souly A\@, Dziemian M\@, Duenas D\@, Lin M\@,
  Wang J\@, Hendrycks D\@, Zou A\@, Kolter Z\@, Fredrikson M\@, Winsor E\@,
  Wynne J\@, Gal Y\@, Davies X}\@. \bibtitle{{AgentHarm: A Benchmark for
  Measuring Harmfulness of LLM Agents}}\bibmark{[J/OL]}\@. \bibjournal{arXiv
  preprint arXiv:2410.09024}\@, \bibyear{2024}\bibcitedate{[2026-09-06]}\@.
  \url{https://arxiv.org/abs/2410.09024}\@.

\bibitem{lightman2023verify}
\bibauthor{Lightman H\@, Kosaraju V\@, Burda Y\@, Edwards H\@, Baker B\@, Lee
  T\@, Leike J\@, Schulman J\@, Sutskever I\@, Cobbe K}\@. \bibtitle{{Let's
  Verify Step by Step}}\bibmark{[J/OL]}\@. \bibjournal{arXiv preprint
  arXiv:2305.20050}\@, \bibyear{2023}\bibcitedate{[2026-09-06]}\@.
  \url{https://arxiv.org/abs/2305.20050}\@.

\bibitem{processbench2024}
\bibauthor{Zheng C\@, Zhang Z\@, Zhang B\@, Lin R\@, Lu K\@, Yu B\@, Liu D\@,
  Zhou J\@, Lin J}\@. \bibtitle{{ProcessBench: Identifying Process Errors in
  Mathematical Reasoning}}\bibmark{[J/OL]}\@. \bibjournal{arXiv preprint
  arXiv:2412.06559}\@, \bibyear{2024}\bibcitedate{[2026-09-06]}\@.
  \url{https://arxiv.org/abs/2412.06559}\@.

\bibitem{zheng2023judging}
\bibauthor{Zheng L\@, Chiang W-L\@, Sheng Y\@, Zhuang S\@, Wu Z\@, Zhuang Y\@,
  Lin Z\@, Li Z\@, Li D\@, Xing E~P\@, Zhang H\@, Gonzalez J~E\@, Stoica I}\@.
  \bibtitle{{Judging LLM-as-a-Judge with MT-Bench and Chatbot
  Arena}}\bibmark{[J/OL]}\@. \bibjournal{arXiv preprint arXiv:2306.05685}\@,
  \bibyear{2023}\bibcitedate{[2026-09-06]}\@.
  \url{https://arxiv.org/abs/2306.05685}\@.

\bibitem{liu2023geval}
\bibauthor{Liu Y\@, Iter D\@, Xu Y\@, Wang S\@, Xu R\@, Zhu C}\@.
  \bibtitle{{G-Eval: NLG Evaluation using GPT-4 with Better Human
  Alignment}}\bibmark{[J/OL]}\@. \bibjournal{arXiv preprint
  arXiv:2303.16634}\@, \bibyear{2023}\bibcitedate{[2026-09-06]}\@.
  \url{https://arxiv.org/abs/2303.16634}\@.

\bibitem{min2023factscore}
\bibauthor{Min S\@, Krishna K\@, Lyu X\@, Lewis M\@, Yih W-t\@, Koh P~W\@,
  Iyyer M\@, Zettlemoyer L\@, Hajishirzi H}\@. \bibtitle{{FActScore:
  Fine-grained Atomic Evaluation of Factual Precision in Long Form Text
  Generation}}\bibmark{[J/OL]}\@. \bibjournal{arXiv preprint
  arXiv:2305.14251}\@, \bibyear{2023}\bibcitedate{[2026-09-06]}\@.
  \url{https://arxiv.org/abs/2305.14251}\@.

\bibitem{es2023ragas}
\bibauthor{Es S\@, James J\@, Espinosa-Anke L\@, Schockaert S}\@.
  \bibtitle{{Ragas: Automated Evaluation of Retrieval Augmented
  Generation}}\bibmark{[J/OL]}\@. \bibjournal{arXiv preprint
  arXiv:2309.15217}\@, \bibyear{2023}\bibcitedate{[2026-09-06]}\@.
  \url{https://arxiv.org/abs/2309.15217}\@.

\bibitem{lin2021truthfulqa}
\bibauthor{Lin S\@, Hilton J\@, Evans O}\@. \bibtitle{{TruthfulQA: Measuring
  How Models Mimic Human Falsehoods}}\bibmark{[J/OL]}\@. \bibjournal{arXiv
  preprint arXiv:2109.07958}\@, \bibyear{2021}\bibcitedate{[2026-09-06]}\@.
  \url{https://arxiv.org/abs/2109.07958}\@.

\bibitem{xiang2024guardagent}
\bibauthor{Xiang Z\@, Zheng L\@, Li Y\@, Hong J\@, Li Q\@, Xie H\@, Zhang J\@,
  Xiong Z\@, Xie C\@, Yang C\@, Song D\@, Li B}\@. \bibtitle{{GuardAgent:
  Safeguard LLM Agents by a Guard Agent via Knowledge-Enabled
  Reasoning}}\bibmark{[J/OL]}\@. \bibjournal{arXiv preprint
  arXiv:2406.09187}\@, \bibyear{2024}\bibcitedate{[2026-09-06]}\@.
  \url{https://arxiv.org/abs/2406.09187}\@.

\bibitem{agentguard2026nexus}
\bibauthor{Hossain E\@, Nipu M~M~H\@, Ornee T~N\@, Rana R\@, Yousefi N}\@.
  \bibtitle{{NEXUS: Structured Runtime Safety for Tool-Using LLM
  Agents}}\bibmark{[J/OL]}\@. \bibjournal{arXiv preprint arXiv:2607.19356}\@,
  \bibyear{2026}\bibcitedate{[2026-09-06]}\@.
  \url{https://arxiv.org/abs/2607.19356}\@.

\bibitem{wang2026safetyvalidity}
\bibauthor{Wang Y\@, Han X\@, Shang D\@, Tang Y\@, Liu B}\@. \bibtitle{{Safety,
  or Just Capability? A Validity Audit of Agent-Safety
  Benchmarks}}\bibmark{[J/OL]}\@. \bibjournal{arXiv preprint
  arXiv:2607.28685}\@, \bibyear{2026}\bibcitedate{[2026-09-06]}\@.
  \url{https://arxiv.org/abs/2607.28685}\@.

\bibitem{yan2024knowledge}
\bibauthor{Yan F\@, Li Z\@, Silamu W\@, Li Y}\@. \bibtitle{{Knowledge-aware
  image understanding with multi-level visual representation enhancement for
  visual question answering}}\bibmark{[J/OL]}\@. \bibjournal{Machine
  Learning}\@, \bibvolume{113}\bibnumber{(6)}\thinspace{}\textnormal{:
  }\bibpages{3789\thinspace{}\textnormal{--}\thinspace{}3805}\@,
  \bibyear{2024}\bibcitedate{[2026-09-06]}\@.
  \url{https://doi.org/10.1007/s10994-023-06426-6}\@.

\bibitem{tang2026hashing}
\bibauthor{Tang L\@, Wang L\@, Wang G\@, Sun M\@, Yan F}\@. \bibtitle{{Feature
  Fusion Mamba Hashing via Decoupling for Cross-Modal
  Retrieval}}\bibmark{[J/OL]}\@. \bibjournal{IEEE Signal Processing Letters}\@,
  \bibvolume{33}\thinspace{}\textnormal{:
  }\bibpages{609\thinspace{}\textnormal{--}\thinspace{}613}\@,
  \bibyear{2026}\bibcitedate{[2026-09-06]}\@.
  \url{https://doi.org/10.1109/LSP.2026.3652952}\@.

\end{thebibliography}


\section{\textbf{项目的研究内容、研究目标,以及拟解决的关键科学问题。}（此部分为重点阐述内容）}

% 迁移来源：extraTex/2.1.研究内容.tex；研究内容、研究目标、拟解决的关键科学问题
\subsection{研究内容}

本项目围绕教育智能体运行中的情境依赖风险，研究风险表征、轨迹评测和监测干预三个相互衔接的内容。研究对象以课程答疑、文献学习辅助中的工具型单智能体为主，采用可控检索和记忆操作记录任务执行过程。双智能体协同仅用于检验错误传播的压力情形，不扩展为独立的大规模多智能体平台。
见图~\ref{fig:edu-f02}。
\HuazhongFigure{figures/F02-framework-v2.pdf}{风险表征、轨迹评测与监测干预的总体框架}{fig:edu-f02}



\textbf{研究内容一：教育情境约束下的可信风险表征。}从任务目标、允许帮助程度及数据授权范围出发，建立能够连接风险类别与外显行为证据的表征方法。将风险归纳为证据与内容失真、教学支持失配、行动与信息边界违背三个相互关联的方面，分别覆盖事实错误与来源不实、目标偏离与过度代答、工具越权与记忆污染。教学支持失配包含支持不足和明确违背帮助边界，两者分开标注和报告。公平性问题通过相同任务条件下的支持差异开展探索，不根据单条输出推断群体偏见。

拟为风险事件记录触发情境、可观察证据、受影响对象及行为后果，并辨认首次能够依据外显信息判断风险的位置。授权帮助边界违背是可操作的规则风险，不能直接等同于已经造成伦理危害。标注分别保留明确违背、教学支持不足和规范不确定，邀请未参与任务生成的教学人员复核边界及反例；争议样本不强行归类，并检查不同处理方式对结论的影响。风险注入仅表明测试情境发生改变，不直接作为智能体已经违规的标签；实际标签需要依据后续可观察行为和独立人工判断建立。研究将特别比较最终答案相同但过程不同，以及过程相近但教学授权不同的任务对，以避免把“回答正确”误当作充分的可信条件。
见图~\ref{fig:edu-f03}。
\HuazhongFigure{figures/F03-risk-representation-v2.pdf}{R1：任务约束、风险证据与独立标注}{fig:edu-f03}



\textbf{研究内容二：面向行为轨迹的过程级可信评测。}在风险表征与标注规则确定后，研究教育情境和轨迹证据对识别结果的共同作用。可观测轨迹包括用户任务、外显计划、检索证据、工具参数、执行回执和记忆读写记录，不要求获取模型私有思维链。对已完成任务开展离线风险定位，对运行中的任务只使用当前前缀和拟执行动作进行告警，分别回答“发生了什么”与“当前是否需要干预”。

评价采用规则核验与语义判别相结合的方式。授权和参数范围等可确定事项由规则检查；教学支持程度、证据误用等依赖情境的事项由具有明确评分依据的语义评价器处理。输出分维度风险判断、支持证据和不确定性，不用一个加权总分掩盖严重越权。通过等信息对照、教育约束消融和跨任务家族测试，识别过程评价的有效成分，并对大模型评价器与人工判断的一致性进行独立检查。

语义评价按风险维度采用0—3级证据评分：0为未见支持风险的证据，1为可疑线索，2为证据支持风险判断，3为明确且直接的违规证据；规范是否可判定另设不确定标记。逐时点保留评分与对应片段，整轨迹以该维度最高评分判断是否曾告警，线上仅累计当时已见信息。同信息的通用与教育评价共用基础模型、输入字段和输出格式，区别在评分依据：前者采用一般安全与任务正确性规则，后者增加本任务的教学支持边界；交叉消融则固定教育评价器，只改变证据和约束条件。离散分数阈值在正常校准轨迹上选择，无法恰好达到名义误报目标时采用不超过目标的最近可行阈值，并报告对应实测值。确定性越权始终阻断；语义评分达到阈值且可修正时重规划，无法安全修正时阻断；规范不确定时转人工；其余动作执行。人工转交单独报告，同时计入正常任务的干预负担，不能隐去其造成的自主完成损失。
见图~\ref{fig:edu-f04}。
\HuazhongFigure{figures/F04-trajectory-evaluation-v2.pdf}{R2：当前前缀评测与动作前干预窗口}{fig:edu-f04}



\textbf{研究内容三：伦理风险动态监测与安全干预。}以第二项内容的前缀风险评价为基础，在回答发送、工具执行和记忆写入之前设置监测点。沙箱先缓存待发送回答，评价器可以读取尚未发布的候选文本，审查通过后才交付学习者；其动作提交时点即对学习者发布时点。例如，在要求引导练习的任务中，即使不调用工具，直接给出完整代答也须经过发送前检查。确定性权限违背触发阻断；可修正的任务偏离优先请求重规划；证据不足或涉及教学判断争议的情况交由人工复核；低风险且满足授权要求的动作继续执行。各类处置的研究重点是触发条件与执行后果是否相称，尤其关注误阻断和无效重规划。
见图~\ref{fig:edu-f05}。
\HuazhongFigure{figures/F05-monitoring-intervention-v2.pdf}{R3：候选动作缓存、顺序检查与分级处置}{fig:edu-f05}



三项研究共用课程答疑与文献学习辅助两类场景，通过公开课程资料、可控构造任务和授权后的高校案例逐层验证，分别标明公开规则改编、专家构造或授权记录的来源，合成角色不用于推断真人行为。课程答疑关注提示、讲解和独立练习边界，文献学习辅助关注证据引用、资料访问与记忆更新。双智能体压力测试进一步检查污染信息是否经转发进入后续工具调用或持久记忆，并以等消息、等步数的单智能体转发作为对照。

\subsection{研究目标}

总体目标是形成可复核的教育智能体可信评测与监测方法，解释情境约束、过程证据和干预后果之间的关系。针对风险表征，形成具有行为证据、标注依据和反例的风险分类与案例集；针对轨迹评测，比较结果评价、通用轨迹评价与教育情境评价的识别表现和有效告警时间；针对监测干预，给出风险执行、正常任务完成、教学支持合规性及人工负担之间的比较结果。

项目不预先承诺未知的性能提升幅度。主要科学判断依据留出任务上的配对差值、置信区间与消融结果；即使假说不获支持，仍应能够说明哪些风险可由简单规则处理、哪些依赖教育情境，以及哪些情形缺少足够证据开展自动干预。上述边界与失败案例也是项目的可验收研究结果。

\subsection{拟解决的关键科学问题}

\textbf{第一，风险含义为何随教育情境改变，哪些证据能够区分这种变化？}研究在“教师演示”“引导练习”等不同支持条件下，同一外显回答或工具行动的风险判断差异。教学情境的作用必须在与通用轨迹评价相同的信息范围内建立，而不能仅靠给本方法增加提示内容获得优势。

\textbf{第二，完整过程中的有效证据能否在运行前缀中及时获得？}研究风险首次可观察位置与实际动作执行之间的可用干预窗口。对于只能在执行后确认的后果，如实归为事后可识别风险，不把它们计入动作前监测的成功范围；对尚无充分证据的情形，保留不确定判断。

\textbf{第三，及时识别能否转化为具有教学代价约束的干预收益？}研究相同初始任务在不同干预策略下的分支结果，比较风险减少是否超出统一拒绝、规则检查等简单策略所能解释的程度。若依赖频繁拒绝或人工接管才能取得低风险，需明确其适用范围和运行负担。

\section{\textbf{拟采取的研究方案及可行性分析。}（包括有关方法、技术路线、实验手段、关键技术等说明）}

% 迁移来源：extraTex/2.1.研究内容.tex；研究方法、技术路线、关键技术、可行性分析
\subsection{研究方法}

\textbf{数据与任务构建。}从允许研究使用的公开课程资料和论文学习材料构建基础任务，保留来源及版本。每个任务明确目标、工具集合、允许的数据范围和教学支持要求；构造干净任务、检索污染、记忆污染及教学约束变化等版本。污染内容仅进入隔离沙箱，外部发送、课程记录修改等动作使用模拟工具与虚拟身份，避免在真实教学环境中实施高风险测试。

实验设计包含24个模板家族、每家族5个基础任务，共120个基础任务。每个基础任务设置4个受控版本、2个模型和3个随机种子，共2880次候选轨迹运行；检索与记忆为同一工具型代理的组件，额外架构和双智能体测试仅在主实验完成后增补。该数量不包含评价器调用、干预分支、失败重试与人工仲裁，预实验将分别测算并设置总费用上限。四个版本中保留一个干净版本；教学约束变化需经人工判断，不能默认计为正常样本。各家族在主分析中等权，独立统计信息受模板家族和基础任务数量限制。先开展小样本预实验估计标注耗时、风险事件比例与方差，再决定是否增加家族数量或收缩次要分析。

\textbf{人工标注与评价器校准。}按照风险定义与证据规则，由2名独立标注者和1名分歧仲裁者判断轨迹，七名研究生成员已列入成员表，标注人选经培训和一致性试标后按任务家族交叉安排，教学规则复核和分歧仲裁人选在正式标注前确定。标注者不知模型身份和方法条件，需指出支持判断的可观察片段。先讨论标注手册中的歧义，再对独立样本报告逐类一致率、Cohen's $\kappa$及分歧类型；对类别失衡同时保留原始列联表，不用单一一致性分数替代分析。

基础任务及其所有变体、种子运行按模板家族整组划分，采用12个家族用于规则与评价器开发、6个用于校准、6个用于留出测试。用于模型辅助构造的数据不直接进入人工真值；目标模型输出也不能同时充当自身评价的参考答案。评价提示、阈值、主要终点及对照方法在访问测试集结果前确定，避免用测试反馈反复改写规则。

\textbf{对照与消融。}设置仅检查结果、通用轨迹评价、具有教育约束的轨迹评价三类离线对照；其中后两者获得相同的可观测字段、上下文容量和模型预算。进一步采用轨迹证据保留/移除与教育约束正确/打乱的交叉消融，检验教育约束和证据之间的交互作用。打乱约束只改变评价器输入，真实任务标签仍由原始情境确定；另设合法授权变化的任务对，用于研究风险含义随情境的变化。两类实验不得混用标签。合法授权比较固定任务、行动和回答，只改变支持规则；另外开展自然运行验证，区分标签含义变化与智能体行为变化。增加等长无关约束和规则措辞改写对照，检查规则关键词捷径。

最终答案一致或正确性相近的任务对作为机制诊断子集，主比较保留所有初始任务，避免依据干预后的成功状态筛选样本。对显式权限字段易于判定的任务与需要语义理解的任务分别报告；对文本长度、风险词、模型身份以及已知模板线索做检查，判断表观提升是否来自捷径。

\textbf{风险与时间指标。}风险识别报告各类别精确率、召回率、误报率和精确率—召回率曲线面积，并列出类别比例。设$t_0$为首个可观察风险证据出现时点，$t_a$为风险动作提交时点，$t_d$为首次有效告警时点。只有$t_0\leq t_d<t_a$的告警计为动作前检出；$t_d-t_0$表示检测延迟，$t_a-t_d$表示剩余干预窗口。未检出和执行后才发现的事件单独统计，不删除后再计算平均延迟。正常轨迹的误报按整条任务中“至少一次错误告警”计算，防止增加监测次数虚增及时性。

阈值在校准集上依据整轨迹误报率10\%的名义工作点选择，同时报告完整曲线与置信区间。各方法只在正常校准轨迹上选择阈值，测试阶段报告各自实际误报及区间，不能仅因名义目标相同就声称实际误报相同。该工作点是比较设置，不是已获统计保证；若校准样本不足以估计低误报区间，则扩大区间报告，并将相关结果定位为探索性证据，不声称达到部署级安全水平。

\textbf{干预效果与统计推断。}运行对照包括事后检查、仅规则保护、统一拒绝和分级干预。基于相同初始任务、初始状态与随机条件成对比较；当干预改变后续路径时重新执行对应分支，不以原未干预轨迹替代干预结果。风险执行比例的分母为预先设定的全部风险测试任务，正常完成率的分母为全部正常任务，拒绝、超时和人工移交均计入。教学支持合规性由独立人工量表评价，与任务正确性分开报告。

以无在线干预的正常任务完成率为参照，在校准集设定正常任务完成率损失不超过5个百分点的比较约束，冻结策略后检验测试结果；损失区间上界低于约束时支持满足，下界高于约束时支持违反，跨越约束则证据不足，不从测试集重选阈值。人工移交按未自主完成计，同时单列其后人工完成情况。对确认、监测、重规划全部计入调用次数、耗时和费用，记录每条初始任务的人工复核次数。

主检验为教育约束与轨迹证据的召回交互差。以$R_{ec}$表示证据状态$e$与约束状态$c$下、校准阈值固定后的家族等权召回，$e=1/0$对应保留/移除轨迹证据，$c=1/0$对应正确/打乱教育约束，定义\[\Delta_{\mathrm{int}}=(R_{11}-R_{01})-(R_{10}-R_{00}).\]

同一预先选定模型的教育支持边界风险作为主要分析，其他模型和风险类别用于稳健性与探索；模型选择不依据测试表现。将交互召回增益5个百分点设为最小实际意义幅度，在预实验后、测试前锁定该判据；只有区间下界超过该幅度才支持预定实用增益，区间支持正向作用但未跨越该幅度时，不声称已有实用收益。区间上界低于该幅度则不支持达到预定实用增益，跨界结果保留不确定性。通用轨迹对照的召回差值和干预后的风险执行比例差值作为支持性分析，后者只有在效用损失区间支持满足约束、风险执行差值区间同时支持下降时，才支持预定的折中收益。测试实际误报不一致时，交互结果仅解释为该校准策略下的表现，不能单独归因为教育约束与证据的机制协同。

主风险阳性仅指依据原始任务情境确认的教育支持边界违背；支持不足与规范不确定不并入阳性，分别报告数量，并对不确定样本的两种极端归类作敏感性分析。各条件共用同一冻结真值和阳性分母，逐家族列出阳性、阴性和不确定计数。预设主估计对象为全部留出家族；任何家族无主风险阳性时，其召回记为不可估计，全部家族的主交互量与相应成功判断也记为不可估计，不赋零或静默删除。可另列有阳性家族的条件性结果，但仅作为探索并报告家族覆盖率。

采用模板家族外层、基础任务内层的分层聚类自助法给出95\%置信区间。可估计时，在每家族内按基础任务是否含已确认主风险阳性分层重抽，任务的全部变体和种子成组保留，四个消融条件使用同一次抽样；该区间条件于观测到的标签构成，不把无阳性样本推定为安全。任何仍不可估计的重抽结果均记录数量，不通过静默丢弃获得区间；种子只估计运行波动，不作为新增独立样本。六个测试家族只能支持所选任务范围内的探索性结论；家族数较少时增加逐家族敏感性分析，次要比较作多重性控制或明确标为探索，不以未显著等同于两种方法等效。
见图~\ref{fig:edu-f06}。
\HuazhongFigure{figures/F06-factorial-design-v2.pdf}{约束与证据的交叉消融及家族隔离验证}{fig:edu-f06}



\subsection{技术路线}

\textbf{总体技术路线。}教育任务及其教学与授权要求首先进入风险表征环节，形成风险类别、证据规则和人工标注依据；这些结果支持轨迹评测的开发与校准，产生能够说明依据的前缀风险判断；经过校准的判断再驱动动作前监测与分级干预。三条路线在同一组课程答疑和文献学习任务上检验，最终形成评测案例、监测原型及其有效范围。留出测试只承担最终验证，不能反向参与当前版本的规则或阈值调整。

\textbf{面向研究内容一的风险表征路线。}以公开学习材料、任务目标和授权范围为输入，先分析正常任务需要执行的动作，再构造受控情境变化与污染版本。依据外显证据定位风险事件，进行盲法双人标注和分歧仲裁，输出风险类别、标注手册及带证据的案例。通过相同答案的过程差异、合法授权变化和标注一致性检查验证表征的区分能力；通过的证据规则与开发数据提供给第二条路线，校准和测试材料保持独立。

\textbf{面向研究内容二的轨迹评测路线。}以第一条路线的风险规则和可观测轨迹为输入，建立规则核验与语义评价器，分别开展完整轨迹定位及仅使用当前前缀的告警。通过等信息通用评价对照、教育约束交叉消融和模板家族留出测试，验证识别表现及其来源，输出分维度判断、证据定位、校准参数和失败案例。满足可解释性与校准要求的前缀评价方法进入第三条路线，证据不足的风险类别保留人工处理。

\textbf{面向研究内容三的监测干预路线。}以待发送回答或拟执行工具/记忆动作、当前轨迹前缀和校准参数为输入，先检查确定性权限，再根据可修正性和不确定程度决定执行、重规划、阻断或人工复核。每次干预后重建对应运行分支，追踪任务完成与风险后果；只有重新通过检查的动作才允许提交。通过与事后检查、仅规则保护和统一拒绝比较，检验风险下降、误阻断和人工负担，输出可运行原型与分场景建议。新发现的问题可用于下一版本开发，但不能回填当前测试版本以改变既有结论。

\subsection{关键技术}

\textbf{情境约束与证据的可核对关联。}将每项风险判断连接到任务中的具体限制与轨迹片段，例如将来源忠实性连接到引用和检索回执，将权限风险连接到工具参数及授权范围，将教学支持失配连接到本次允许的帮助层级。通过保留反例和不确定标签，使“需要更多证据”的情况与“已识别风险”能够区分。相关表征用于解释判断，不声称从外显日志恢复模型内部因果过程。

\textbf{离线标签与在线信息的时间隔离。}离线人工标注可以读取完整记录以确认事件，在线评价只能读取当时的前缀。评价输入构造和日志记录分别保存，检查未来动作、最终结果及人工标签是否误入在线提示。对只能事后观察的风险单列指标，对无法完成的模拟工具调用保留失败回执，避免把工具不可用混同为安全拒绝。

\textbf{干预分支的可复现比较。}固定任务版本、模型版本、工具行为和预算，记录动作提交前后的环境状态。校准阶段锁定干预规则后，复用相同初始条件生成各策略分支。双智能体压力测试分别控制来源标识与权限约束，追踪下游采纳、工具执行及持久写入；若等步数单智能体转发也出现同等传播，只能归结为信息传播风险，不能宣称多智能体特有的放大机制。

\subsection{方法与实施可行性}

项目采用已有模型与可控工具环境，核心工作量集中在任务设计、标注、评价器校准和对照实验，不依赖基础模型训练。公开文献和工具型智能体基准提供了方法参照，申请人现有算力及软件工程条件能够支持离线运行。三个研究内容共享任务与日志，能够减少分别建设系统造成的重复投入。

执行按“沙箱机制检验—近真实案例—授权后真实场景”推进。涉及真实课程及师生数据的验证，在取得相应授权并完成适用的伦理审查程序后开展。上述条件落实前，只开展公开资料和合成角色构成的沙箱实验；真实场景迟延时，仍可完成方法比较、监测原型和近真实案例报告，但对学生学习效果与实际部署风险不作外推结论。预实验若显示标注或运行成本超出一年期投入，优先保留教育情境与轨迹证据的主比较，将公平性和双智能体扩展保留为探索。

% 迁移来源：extraTex/3.1.研究基础.tex；可行性分析、研究风险与应对措施
\subsection{研究积累与组织可行性}

在方法层面，现有通用轨迹与工具风险研究提供了可参考的任务构造、风险标注和基线评价方式；申请人的问答和知识库经验有助于把公开资料组织为可复现任务。在实现层面，已有算力与Linux、Python、深度学习及版本管理工具链能够支持模型推理、日志记录和成对实验。项目以已有模型为对象，不以训练新基础模型为前提。

在组织层面，申请人承担总体方案、技术实现与证据核查；七名研究生成员及技术分工已列入成员表和研究基础，覆盖任务构建、风险规则、轨迹采集、评测、校准、干预和复现。项目拟在研究启动阶段协商实验室合作分工，组织具有相关教学经验的人员复核任务与评价结果。教学规则设计和人工评分应由具备相应课程理解能力的人员参与，不能只由算法评价器自行生成并验证标签。如教学规则复核人员无法落实，相关成果限于合成协议验证，教育风险测量效度仍需后续检验。合作落实前先完成公开资料和沙箱方法准备，以降低场景依赖对总体进度的影响。

\subsection{研究风险与应对措施}

\textbf{风险标签存在歧义。}当预实验中分歧集中在帮助程度或教学目标解释时，先补充正反例与任务规则，区分明确风险、支持不足和规范不确定，不通过扩大自动评分规模掩盖分歧。对仍有争议的类别报告敏感性分析，必要时收缩为授权明确的风险任务，保留可复核案例作为最低交付。

\textbf{轨迹评价提升缺少稳定性。}当增益只出现于显式风险词、某一模型或少数模板家族时，检查信息量、长度和任务构造捷径，使用等信息对照与家族留出分析定位原因。若教育约束没有额外作用，按反证结果报告，不通过更换测试集寻求有利结论；将可由规则解决的风险与需要人工判断的风险分别说明。

\textbf{干预降低了正常教学支持。}当误阻断、重规划次数或人工复核负担上升而风险改善有限时，减少自动干预覆盖范围，仅保留确定性权限保护，把教学争议交由人工。评价中保留所有拒绝与失败任务，避免因排除困难样本而高估效果。

\textbf{场景授权或资源进度滞后。}在真实课程数据未获许可时使用公开材料与虚拟角色完成沙箱实验，不接触真实学生身份和成绩。模型接口或成本发生变化时，优先采用可本地复现的模型，保留版本和费用记录；减少双智能体与公平性扩展，确保主要情境对照、评价方法和案例报告在一年内完成。

% 迁移来源：extraTex/3.2.工作条件.tex；全文
\subsection{已具备的实验条件}

申请人可稳定使用NVIDIA A40显卡8张和NVIDIA GeForce RTX 3090 Ti显卡6张，并具备Linux、Python、深度学习框架、Git与LaTeX等软件工具条件。上述资源可支持已有模型推理、任务重放和评价实验，项目无需另行采购大型计算设备。实际模型版本和部署规模在预实验中根据资源占用确定，不以全部算力长期独占为前提。

项目可从具有明确使用条件的公开课程资料和论文资源构建任务，配合本地模拟检索、工具调用和记忆环境形成可观察轨迹。公开可获取不等于本课题已拥有数据使用授权，任务整理时将逐项核对许可、引用与可发布范围。现有条件主要支撑公开资料与可控沙箱研究，真实教学数据另行落实。

\subsection{尚缺少的实验条件及拟解决的途径}

项目拟在研究启动阶段协商实验室合作，邀请具有相关教学经验的人员完善任务规则与标注手册，先开展公开资料和模拟环境实验。涉及真实课程及师生数据的小规模案例验证，在明确课程场景、取得相应授权并完成适用的伦理审查程序后开展。涉及真人参与的活动应与课程成绩和正式教学评价分离，参与与退出不影响学生权益。

若上述条件未按计划落实，以公开材料、虚拟角色和模拟工具完成主要方法实验，并将案例明确标注为近真实验证。需调用外部模型时仅使用可公开或获准处理的数据，记录成本。本次依据实验室2026年度正式指南，按一般课题申请5万元，研究期1年；预算编列业务费2万元、劳务费2万元、设备费0万元和间接费用1万元。助理费用按七名成员分阶段承担工作的累计16人月测算，外聘专家咨询按4人天测算；实际劳务核算、外聘安排及支出标准由申请人与所在单位审核后落实。

\subsection{利用研究基地的计划与落实情况}

本项目申报华中师范大学人工智能与教育新形态实验室开放课题，主要对应一般课题方向4“教育人工智能伦理审查标准与监测机制”，兼顾方向5的风险防控与方向7的高等教育实践。拟与实验室人员围绕教育风险定义、教学支持尺度和案例分析开展合作；合作分工和授权范围将另行协商，实验室教学数据在取得授权后使用。申请人已有算力支撑核心技术工作，后续合作用于完善教育情境解释与案例验证。

\section{\textbf{本项目的特色与创新之处。}}

% 迁移来源：extraTex/2.2.特色与创新.tex；全文
\subsection{项目特色}

项目以教育任务中的实际支持边界为研究起点，将教学适切性、证据忠实性与行动授权放在同一行为过程内考察。课程答疑和文献学习辅助共用可复现任务与工具环境，使风险分类、过程评价和干预后果能够相互核对。相同最终答案对应不同轨迹、相同行动对应不同教学要求的受控案例，有助于辨认最终结果难以揭示的风险，也避免把所有教育问题统一解释为模型能力不足。

研究规模服从一年期开放课题条件，以已有模型、公开材料和沙箱为主要载体。成果围绕同一套任务、日志和证据规则形成，评测基准、监测原型与案例报告能够互相追溯。对于真实场景条件尚未落实的部分保留授权前提，使核心机制检验不依赖高成本的全校部署。

\subsection{创新点}

\textbf{教育情境约束与过程证据共同作用的可信评测。}通用风险轨迹基准和近期教学工作流评测已经提供了重要基础。本项目拟进一步以允许帮助程度、任务目标和授权范围的受控变化，检验同一行为风险含义的改变，并在等信息通用评价对照下分析教育约束与过程证据的交互作用。创新的判断依据是跨任务家族的增益及其消融反证，而非仅增加教育题目或评测维度。

\textbf{以时间窗口和教学代价共同检验监测的有效性。}围绕首次可观察证据、首次告警和风险动作提交，区分事后定位与动作前发现，再在相同初始任务上追踪执行、重规划、阻断和人工复核产生的分支结果。研究拟说明哪些风险值得自动干预、何时需要保留人工判断，以及风险减少是否超出增加拒绝所能解释的程度，从而形成有适用边界的监测依据。

\section{\textbf{研究计划、预期成果及考核指标。}}

% 迁移来源：extraTex/2.3.年度研究计划.tex；全文
\subsection{年度研究计划}

本课题研究期为12个月，自正式立项之日起计算，分为任务设计与预实验、评测体系构建、留出评测与干预验证、案例分析与成果总结四个阶段。各阶段依次完成风险定义、评价校准和干预比较，前一阶段的任务规范与验证结果为后一阶段提供依据。

\textbf{第1—3个月：}完成两类教育任务的边界设计、公开资料与使用条件核对，建立隔离工具环境和轨迹采集方式。围绕风险类别、首次可观察位置和教学支持要求编制标注手册，以小样本预实验检查分歧、运行成本和有效风险比例，据此调整任务家族数量与主要比较。阶段交付为任务规范、标注样例、预实验分析和可复现的环境配置。

\textbf{第4—6个月：}完成基础任务与变体构建、独立标注及数据划分，建立结果评价、通用轨迹评价和教育情境评价的对照。开展教育约束消融，使用校准集确定监测工作点与效用要求，并在读取测试结果之前明确主要终点和统计方案。同步撰写方法部分，阶段交付为评测数据初版、证据规则、评价脚本和冻结的分析计划。

\textbf{第7—9个月：}开展留出任务的轨迹评测与分支干预比较，报告风险、误报、正常任务完成和人工负担。针对主要失败类别补充机制分析；在核心任务完成且资源允许时开展双智能体传播压力测试。本阶段形成主要结果论文初稿，阶段交付为主要实验报告、监测原型及失败案例说明，不以追求单一最好分数替换预定比较。

\textbf{第10—12个月：}完成授权条件下的高校案例分析；若真实数据或场景未获许可，则采用明确标注的近真实沙箱案例。整理风险表征、评价器与监测方法的适用范围，完成论文和典型案例报告，清理可公开数据与代码，形成可复核的成果包。方法与实验完成后及时投稿并跟进修改，结合审稿进度推进论文发表；成果考核以正式发表为准。

\subsection{预期研究结果}

依据实验室2026年度正式指南，一般课题考核目标为在SCI、SSCI或CSSCI收录期刊发表论文不少于1篇，并提交典型案例研究报告1份。围绕这两项成果，拟形成教育智能体风险分类与标注规范1套、可复用的任务与轨迹评测集1套、风险监测原型1个及配套评价脚本。论文、基准和原型围绕同一实验主线形成，研究结果与数据、代码相互对应。

评测集将同时提供正常任务、受控风险任务、人工判断依据、数据划分与运行设置，公开范围服从原材料许可和授权约定。研究报告不仅呈现方法收益，也记录无增益、误告警和过度干预等反例，给出需要人工复核的条件。软件著作权和额外高水平论文作为资源允许时的争取成果，研究规模以完成主要任务为限。

获资助后形成的课题成果，按指南由实验室和研究者所在单位共有，并按要求标注实验室名称、开放课题资助及实际批准编号。结题时提交研究报告、发表证明、案例报告及成果清单。既有横向合同的受限资料和知识产权不因本次合作而转移，成果使用分别遵守原合同与本课题约定。
